\cotefile{10}
\scanpage{155}
Soit $G$ un \unsure{groupe|grpe} \deleted{\unsure{quel}} algébrique linéaire quelconque,
$T$ un tore maximal, $N$ son normalisateur et $W=N/T$.
\deleted{\unsure{G}} $N$ opère sur $G$ et $T$ de façon compatible par
\unsure{autom.} intérieurs, et \deleted{\illegible{1}} donc \deleted{\unsure{le}} l'hom canonique
\[ CK(G)\longrightarrow CK(T)\simeq\mathbb Z[\widehat T] \]
est compatible avec les opérations de $N$. $N$ opère
trivialement sur le pr. mb., et $T$ opère trivialement
sur le second, donc en fait $W$ opère sur
$\mathbb Z[\widehat T]$ via les opérations de $W$ sur $\widehat T$. On
a donc un hom
\[ CK(G)\longrightarrow\mathbb Z[\widehat T]^W \tag*{(5)} \]
où $\mathbb Z[\widehat T]^W$ désigne \deleted{le}l'\deleted{groupe}\inserted{anneau} des \unsure{élémts} de
$\mathbb Z[\widehat T]$ invariants sous $W$.

\deleted{On} En général \deleted{\unsure{(à cause de la torsion)}} on n'a
pas là une bijection. Mais c'est \inserted{sans doute} une bijection
s'il n'y a pas de torsion, \deleted{\unsure{(du moins on}} et de
plus l'homomorphisme analogue
\[ C\bigl(K(G)\mathbin{\underset{\mathbb Z}{\otimes}}k\bigr)\longrightarrow k[\widehat T]^W \tag*{(5 bis)} \]
(où $k$ est un corps commutatif) est bijectif
si \deleted{\unsure{G} n'a pas de} \inserted{la car. de $k$ est} première à la torsion de $G$.

% archive sheet number: 154
% paper watermark: ARGENTINE TYPEWRITER (partially visible at the edge)
\scanpage{157}
$\unsure{\mathcal K|K|\mathcal C}$ classe abélienne, \deleted{\unsure{foncteur}} dont les éléments sont des
\deleted{\unsure{groupes}} \uline{modules} sur l'anneau commutatif $k$
\margin{\unsure{obtenue|obtenus} par l'\unsure{imposition} de structures supplémentaires sur les $k$-modules}
% slanted
et les homomorphismes des hom. de modules,
la composition des hom étant la composition usuelle.
On suppose le foncteur identique $\mathcal K\longrightarrow C^{k}$ \uline{exact}
et qu'il existe un \uline{ensemble} de $A_i\in\mathcal K$ tels
que $\forall\,V\in\mathcal K$ soit isomorphe à l'un d'eux

Notation $\operatorname{Hom}_{\mathcal K}(V,W)$ si $V,W\in\mathcal K$;

le noyau, image \inserted{etc.} au sens $\mathcal K$ \unsure{corresp. et|correspondent} à noyau
image etc au \unsure{sens} \deleted{\unsure{mod}} ordinaire.

\unsure{Les|Soit|Lemme} $I:\mathcal K\longrightarrow C^k$ le foncteur identique. Les endomorphismes
de ce foncteur s'identifient \deleted{\unsure{aux}}\inserted{\unsure{aux systèmes}}
\deleted{homomorphisme}\inserted{d'endomorphismes} de $k$-modules $(u_i)$, $u_i\in$
$L_k(A_i)$; tels que pour \unsure{tout|$\forall$} hom
$\varphi\in L_{\mathcal K}(A_i,A_j)$, on ait commutativité
des
\diagram{\begin{tikzcd}[ampersand replacement=\&] A_i \arrow[r,"\varphi"] \arrow[d,"u_i"'] \& A_j \arrow[d,"u_j"] \& {} \\ A_i \arrow[r,dash,"\varphi"] \& A_j \& \varphi u_i=u_j\varphi \end{tikzcd}}

Ces endomorphismes forment donc une
sous-classe \deleted{\unsure{bije}}\deleted{\unsure{hom}} de $\prod_i L(A_i)$, fermé
pour la topologie produit, qd les
$L(A_i)$ sont munis de la convergence simple
($A_i$ pris discret)

Donc c'est un \uline{anneau \deleted{topologique}} \inserted{muni d'une topologie}
\inserted{linéaire gauche,} \uline{complet}\unsure{,} \inserted{soit $U$.} Si $V\in C$, alors $V$ est un
module sur $U$, et $\underline U\longrightarrow L_{\unsure{s|3|k}}(\unsure{W|V})$ est continu
De plus, les $\varphi\in\operatorname{Hom}_{\mathcal K}(V,W)$ sont
des $U$-homomorphismes. On obtient

% archive sheet number: 156
% faint left-margin pencil marks: 77
\scanpage{158}
\deleted{\unsure{A}} ainsi un foncteur canonique, noté $\unsure{J|I}$, de
$\unsure{K|\mathcal K}$ dans les $U$-modules ``permis'' (i.e. tels que
$U\longrightarrow L_s(W)$ soit continue). On va donner des
conditions pour que ce soit \deleted{bijectif}\inserted{un isomorphisme}. Supposons les
$A_i$ à engendrement fini sur $k$, alors les $L(A_i)$
sont \uline{discrets}, donc $\underline U$ est un anneau \deleted{\unsure{linéairement|linéaire}}\inserted{topologique},
\inserted{(avec une topologie linéaire)}
et on va se borner à considérer les $\underline U$-modules
permis à engendrement fini sur \deleted{$\unsure{A_i}$} $k$. \deleted{Pour ceci,}
 On \deleted{\unsure{va}} suppose ceci\unsure{:}

\deleted{pour tout $V\in K$, \deleted{il y a un \unsure{hom}} soit $\underline U_V$ le commutant de $\operatorname{Hom}_K(V,V)$ dans $\operatorname{Hom}_k(V,V)$. On suppose que tte représentation de $\underline U_V$ dans un $k$-module de type fini $W$ définit sur $W$ \inserted{fonctoriellement} une structure $K$, de façon \unsure{précise} on suppose donné un foncteur \inserted{``permis''} $F_U$ de $C^{\underline U_{\unsure{V|W}}}$ dans $K$ \deleted{qui, une} (i.e. foncteur qui composé avec $I$, donne le \deleted{in\unsure{itial}} foncteur identique $C^{\underline U_V}\longrightarrow C^k$); et \deleted{que tout $K$-hom $V\longrightarrow W$}}

\deleted{\unsure{de} Il existe un foncteur \deleted{\unsure{de}} de $C'^{\underline U}$ (modules permis sur $\underline U$, à \unsure{engendrement} fini sur $k$) dans $\unsure{\mathcal K|K|\mathcal C}$, permis (i.e. \deleted{\unsure{comp}} composé avec $I$, on obtient l'injection $\overline I$ de $C'^{\underline U}$ dans $C'^{k}$) tel \deleted{que} \unsure{le|les}}

% archive sheet number: 157
\scanpage{159}
Pour tout $\unsure{V|W}\in C'^{k}$ ($k$-module de type fini)
soient $K(W)$ \deleted{\unsure{resp.|et}} $K'(W)$ l'ens des ``structures
de type $K$'' resp$^{\mathrm t}$ des structures de $\underline U$-modules
permis sur $k$. On a défini un hom
fonctoriel $K(V)\longrightarrow K'(W)$. \unsure{Ceci|Ceci étant} dit

$1^\circ)$ On veut que cet hom. soit bijectif

$2^\circ)$ Si $V,W\in K$, on veut que les
$K$-hom de $V$ dans $W$ soient identiques aux
$\underline U$-homomorphismes.

On est ramené à la considération des $\underline U$-\deleted{\unsure{lin}}\inserted{modules}
permis qui sont à engendrement fini \deleted{sur} $k$.

$\underline U$ est un \uline{anneau topologique} \inserted{\uline{linéaire, complet avec $1$}\deleted{\unsure{,0}}} ayant un
système \deleted{\unsure{compl}}\inserted{fond.} de \deleted{\unsure{représentations}} \inserted{voisinages $\mathcal J$ \unsure{sts|tels} que}
\deleted{\unsure{de $k$} de $U$ dans $L(V)$, $V$ de type fini sur $k$.} \inserted{$\underline U/\mathcal J$ soit de dim finie.} Tout
anneau de ce type peut s'obtenir ainsi.
\margin{Si $k$ est un corps les $\underline U$ sont les \unsure{duals|duels|duaux} des co-algèbres \unsure{avec} \deleted{\unsure{t}} qui sont réunions \deleted{\unsure{filtrantes de ses}} \deleted{\unsure{algèbres de dim}} de ses co-algèbres de dim finie. Aux hom. continus des $\underline V\to\underline U$ compatibles avec l'augmentation correspondent les hom de Co-Algèbres compatibles avec l'augmentation.}
% slanted

\uline{Première structure supplémentaire}\,: On suppose
\deleted{\unsure{que $k$}} \inserted{tout $V$ \uline{libre} \deleted{\unsure{de $k$}}} puisse être muni canoniquement
d'une structure $\unsure{\mathcal K|K|k}$, i.e. que on ait
sur $k$ une structure privilégiée de $\underline U$ module,
\deleted{\unsure{évidemt}} il revient au $=$ de se donner sur
$\underline U$ une \uline{augmentation continue $\varepsilon$}

% boxed note
\deleted{Les} $\underline U$-modules permis sont les \uline{comodules unitaires}
sur $A$ \deleted{\unsure{qui}} qui sont de dim finie sur $k$.

% archive sheet number: 158
\scanpage{160}
Soient $\underline U,\underline V$ deux \deleted{\illegible{1}} ``$\underline{\unsure{k|K}}$-algèbres'' i.e.
algèbres topologiques \deleted{\unsure{augmentées}} satisfaisant aux
conditions dites plus haut. Un hom. continu
\deleted{compatible avec l'augmentation} $\underline V\longrightarrow\underline U$
définit un \deleted{hom.}\inserted{foncteur} permis $\unsure{\mathscr C|K|\mathcal C}'^{\underline U}\longrightarrow\mathscr C'^{\underline V}$;
\deleted{et on obtient ainsi un fonct et un seul}
\deleted{\unsure{fonct}}\inserted{i.e. foncteur tel que}
\diagram{\begin{tikzcd}[ampersand replacement=\&] \mathscr C'^{\underline U}\arrow[r]\arrow[dr,"I^{\underline U}"] \& \mathscr C'^{\underline V}\arrow[d,"I^{\underline V}"] \\ {} \& \mathscr C'^{k} \end{tikzcd}}
soit commutatif.

Réciproquement, \deleted{\unsure{un}}\inserted{considérons} tel foncteur, \deleted{définit}
\deleted{\deleted{soit $v\in\underline V$.} Soit $V\in C'^{\underline U}$, on va définir une opération $\unsure{\sigma|o}_V$ dans $\underline I(W)$, \deleted{\unsure{on soit}} ceci du fait que on a une opération de $v$ ds $F(V)$ qui est un $V$-module. Les $\unsure{o|\sigma}_V$ définissent \deleted{\unsure{un}} endomorphisme du foncteur}

Comme $I^{\underline U}=I^{\underline V}\circ F$, tout endomorphisme
de $I^{\underline{\unsure{W|V}}}$ en définit un de $I^{\underline U}$, d'où
un hom de l'anneau des end.
de $I^{\underline W}$ dans celui de $I^{\underline V}$, i.e. de $\underline{\unsure{V|W}}$ dans
$\underline U$, dont on vérifie qu'il est continu.

\uline{Ainsi}, \uline{les hom continus} $\underline V\longrightarrow\underline U$
\uline{s'identifient aux hom. permis de foncteurs}.

% archive sheet number: 159
\scanpage{161}
Si on a sur $\underline U$ et $\underline V$ des augmentations continues,
les hom. compatibles avec l'augmentation
correspondent aux foncteurs permis qui transforment
le $\underline U$ module trivial \deleted{\unsure{aug}} $k$ en le $\underline V$-module trivial $\underline k$.
Transitivité, etc.

\deleted{Supposons donné un foncteur}

\uline{Autres structures supplémentaires}

\inserted{Soit $k$ un \deleted{\unsure{fonct}}\inserted{corps}.} Supposons donné un foncteur $F^\vee\ \unsure{\mathcal K|\mathscr C|\mathcal C}^{\underline U}\longrightarrow\mathcal K^{\underline U}$
\deleted{tel que l'on ait une \unsure{identité}}
associant à \deleted{tout} $V\in K^{\underline U}$ une structure
de $\underline U$-module permise sur le dual $V'$ de $V$,
\deleted{De telle façon que}
i.e. un foncteur tel que
\diagram{\begin{tikzcd}[ampersand replacement=\&] \underline K^{\underline U}\arrow[r,"F^\vee"]\arrow[d] \& \underline K^{\underline U}\arrow[d] \& {} \\ \underline K^{k}\arrow[r,"F^\vee"{name=lowerF}] \& \underline K^{k} \& \breve F(W)=V'\arrow[from=lowerF,to=2-3,bend right=15] \end{tikzcd}}
soit commutatif. \deleted{\unsure{Cela comme}}
\deleted{$V'$ est un $\underline U^\circ$-module ($\underline U^\circ$ inverse de $\underline U$) cela signifie}
On veut que de plus,
$F^\vee F^\vee={}\underline{\text{identité}}$. \deleted{Ceci définit, par} \deleted{Il revient au $=$ de}
\deleted{\unsure{se donner}, sur tout} Soit alors $V$ un
\deleted{tout} module permis sur $\underline U^\circ$
(\unsure{l'opposé|l'opposée|l'inverse} de $\underline U$) $V'$ est un module sur
$\underline U$ (transposition) donc $(V')'$ est un
module sur $\underline U$. On a ainsi un
\deleted{\unsure{contrav.}} foncteur permis des $\underline U^\circ$-

% archive sheet number: 160
\scanpage{162}
modules dans les $\underline U$ modules, i.e. un
\deleted{\unsure{isomph|isomorphisme}} homomorphisme $\underline U^\circ\longrightarrow\underline U$, i.e.
un \deleted{\unsure{anti-automorphisme}}\inserted{application $\sigma\colon\underline U\longrightarrow\underline U$ qui est \unsure{$k$-linéaire|$k$-lin.}}
et anti-multiplicative; $\breve F\breve F=I$ signifie
que $\sigma^{\unsure{2}}=1$, donc c'est un \uline{antiautomorphisme involutif}, \uline{noté $x\longrightarrow\breve x$}. Réciproquement,
un tel antiautomorphisme définit
\deleted{\unsure{une}} un foncteur $K^{\underline U}\longrightarrow K^{\underline U}$. Dans l'hyp. ci-dessus, \deleted{\unsure{si|qui}} --- Si on a de
plus une augmentation $\varepsilon$, la ``\uline{compatibilité}
\uline{de $\sigma$ avec $\varepsilon$}, i.e. $\varepsilon(x)=\varepsilon(\breve x)$
signifie que le \deleted{$\underline U$-structure de}\inserted{\uline{dual du $\underline U$ module}}
\deleted{\unsure{du $k$-fonct}}\uline{trivial est le $U$-module trivial}. \deleted{En} Si
on a un \uline{hom} continu $\varphi:\underline V\longrightarrow\underline U$, le
fait qu'il est \uline{compatible avec $\sigma$},
\deleted{Supposons donné un bifoncteur}
i.e. $\varphi(\breve x)=\overset{\vee}{\overline{\varphi(x)}}$, signifie que
\deleted{$F(A,B)$}
\deleted{si} \deleted{$\overline\varphi(V)$} si $V\in K^{\underline U}$, on a
$\overline\varphi(V)'=\overline\varphi(V')$, i.e le foncteur
permis est compatible avec le
passage au dual.

% archive sheet number: 161
\scanpage{163}
Supposons donné un \uline{bifoncteur}
$T(V,W)$\,: \hfill tel que
\diagram{\begin{tikzcd}[ampersand replacement=\&,column sep=small] K^{\underline U}\times K^{\underline U}\arrow[r,"T"]\arrow[d,"I^{\underline U}\otimes I^{\underline U}"] \& K^{\underline U}\arrow[d,"I"] \& {} \\ K^k\times K^k\arrow[r,"T"{name=lowT}] \& k^{k} \& T(\unsure{W|V},\unsure{W|V})=V\underset{k}{\otimes}W \arrow[from=lowT,to=2-3,bend right=20] \end{tikzcd}}
soit commutatif. Cela signifie que si $V,W\in K^{\underline U}$,
on a sur $V\underset{k}{\otimes}W$ une structure de
$\underline U$-module permis, de telle façon que
le produit tensoriel de $\underline U$-hom soit
\deleted{un} $\underline U$-homomorphisme.
\margin{\deleted{On \uuline{suppose} que les opérations de $U$ dans $V\underset{k}{\otimes}W$ sont \unsure{contenues} dans $\underline U_V\otimes\underline U_W$ ($\underline U_V={}$\deleted{\unsure{algèbre}}\unsure{opérateurs} sur $V$ définis par $\underline U$)}}
% slanted

\deleted{Soit \deleted{alors} $X$ un bi-$\underline U$-module permis (i.e. \deleted{$X$ muni} de deux structures \inserted{permis} \deleted{\unsure{commutantes}} de $\underline U$-module, i.e. $X$ est un \deleted{\unsure{bi}} module permis sur $\underline U\underset{k}{\widehat\otimes}\underline U$. Mettons sur $X$ une structure de $\underline U$-module\,: Considérons $\underset{\unsure{\Vert}}{X}\underset{\unsure{\Vert}}{\otimes}X$}

\deleted{Soit alors $V\in K^{\underline U}$, soit $\mathcal J_V$ le noyau de $\underline U\longrightarrow L_k(V)$, et \deleted{on a un hom naturel de $\underline U$ dans $L_k(\unsure{V|W}\otimes V)$ \unsure{Pour $\mathcal J_V$}} $\underline U_V=\underline U/\mathcal J_V$. Tout $\underline U_V$-module $W$ est un $\underline U^{\unsure{\prime}}$ module, \deleted{donc $W\otimes W$ est un $\underline U$-module, i.e. on a une représentation fonctorielle $\underline U\longrightarrow L_k(W)\otimes L_k(W)$} $\underline U_V$ étant un anneau d'Artin, il n'y a qu'un \unsure{nb fini} de $\underline U_V$-modules}

% archive sheet number: 162
\scanpage{164}
\deleted{indécomposables distincts, soient $V_i$. pour $i,j$, $V_i\otimes V_j$ est un $\underline U$-module, donc se décompose en un nb fini de sous $\underline U$-modules \deleted{irréductibles} indécomposables $W_{ij\unsure{k|u}}$, l'ens des $W_{ijk}$ obtenus est fini. \deleted{Soit $\mathcal J'_V$ le noyau} de $U$ dans $\sum W_{ijk}$}

\deleted{Montrons qu'il existe un $W\in\inserted{K^{\underline U}}$ tel que, si $W_1$ et $W_2$ sont $\in K^{\underline U_V}$, alors $W_1\otimes W_2$, muni de sa structure ``naturelle'' de $\underline U$-module, ait un annulateur $\supset\mathcal J_W$. En effet, il suffit de le vérifier pour $W_1$ et $W_2$ indécomposables (\unsure{grâce à} \unsure{Rem th de|Krull} Schmidt) or il n'y a qu'un nb \uline{fini} de $V_i$ indécomposables distincts; on prendra $W=\sum V_i\otimes V_j$. Ceci dit, \deleted{on veut définir} \deleted{\unsure{D'ailleurs}, \illegible{2} un hom}}

Ceci est un cas particulier de la \deleted{\unsure{structure}}\inserted{donnée}
suivante. Tout bi-$\underline U$-module permis
(i.e. un module permis sur $\underline U\mathbin{\widehat\otimes}\underline U$)
définit un $\underline U$-module permis de
façon fonctorielle, i.e. on se donne
un \uuline{hom diagonal} $\underline U\longrightarrow\underline U\mathbin{\widehat\otimes}\underline U$. Mais on

% archive sheet number: 163
\scanpage{165}
va voir que la donnée des \deleted{\unsure{$K$-}}structures \inserted{de $\underline U$ modules} sur les
$V\underset{k}{\otimes}W$ ($V,W\in\underline K^U$) suffit déjà à déterminer
un $\underline U\longrightarrow\underline U\mathbin{\widehat\otimes}\underline U$ \uline{unique} tel que ce soit défini
par lui. Pour ceci, il suffit, pour tout
voisinage $\mathcal J$, de \deleted{\unsure{se}} définir
$\underline U\longrightarrow\underline U/\mathcal J\otimes\underline U/\mathcal J$. \inserted{On peut supposer que} $\mathcal J$ est l'annulateur
d'un $V\in K^{\underline U}$, donc $\underline U/\mathcal J\inserted{={}\underline U_V}\subset L_k(V)$,
$\underline U/\mathcal J\underset{k}{\otimes}\underline U/\mathcal J\subset L_k(V)\underset{k}{\otimes}L_k(W)=L_k W\underset{k}{\otimes}V)$.

\deleted{il suffit donc} Or on \unsure{connaît|connait} bien
$\underline U\longrightarrow L_k(\unsure{V|W}\otimes V)$, cela définit ce qu'on veut
pourvu qu'on fasse la supposition

Les $\underline U$-homothéties de $V\otimes V$
$(V\in K^{\underline U})$ sont de la forme $\sum u_i\otimes v_i$,
où $u_i$ $v_i$ sont des $\underline U$-homothéties
de $V$.
\margin{\deleted{\unsure{Si $K^{\underline U}$ est|$k^{\underline U}\to$}} Si $K^{\underline U}$ est $1/2$ simple, cette condition est vérifiée d'elle-même}
% slanted

Ainsi on définit $\underline U\longrightarrow\underline U/\mathcal J\underset{k}{\otimes}\underline U/\mathcal J$.
On voit aussitôt les compatibilités
nécessaires pour que $\underline U\longrightarrow\underline U\mathbin{\widehat\otimes}\underline U$ soit
\unsure{défini|définié}, et on vérifie que cet hom.
définit la loi tensorielle cherchée.

% archive sheet number: 164
\scanpage{166}
\uline{Axiomes supplémentaires} \uline{Associativité} de
l'homomorphisme $\Delta$\,: correspond à un axiome
d'associativité pour les produits tensoriels.
\margin{On exprime ici que les hom. tensoriels classiques sont des $\underline U$-homomorphismes}
% slanted

\uline{Commutativité} de $\Delta$\,: l'hom canonique
$V\underset{k}{\otimes}W\longrightarrow W\underset{k}{\otimes}V$ doit être compatible
avec les $\underline U$-structures. On ne
parlera par la suite que de bifoncteurs
tensoriels associatifs et commutatifs, i.e. $\Delta$
associatif et commutatif. --- \uline{Compatibilité}
\uline{avec l'augmentation} \inserted{(i.e. $\Delta(u)=1\otimes u+u\otimes1+{}$\unsure{él. mixtes|élém. mixtes} \unsure{si} $\unsure{\varepsilon(u)=0}$)} $\Longleftrightarrow$ Le produit tensoriel
d'un $\underline U$-module $V$ avec le $\underline U$-module
trivial $k$ est le \deleted{$\underline U$-module trivial} $V$

\uline{Compatibilité avec $\vee$}\,: \deleted{\unsure{Si} $k$ $V,W$} Si $V,W\in K^{\underline U}$,
alors $(V\underset{k}{\otimes}W)' \simeq V'\underset{k}{\otimes}W'$ --- \unsure{tant que|en tant que}
$k$-modules, on veut que ce soit vrai
en tant que $\underline U$-modules \unsure{permis|\deleted{permis}}. Ceci
signifie qu'on a commutativité dans
\diagram{\begin{tikzcd}[ampersand replacement=\&] U\arrow[r,"\Delta"]\arrow[d,"\vee"] \& U\otimes U\arrow[d,"\vee\otimes\vee"] \\ U\arrow[r,"\Delta"] \& U\otimes U \end{tikzcd}}
\diagram{\begin{tikzcd}[ampersand replacement=\&] \deleted{\unsure{U}}\arrow[r,dash]\arrow[d] \& \deleted{\unsure{U\otimes U}}\arrow[d] \\ \deleted{\unsure{U}}\arrow[r] \& \deleted{\unsure{U\otimes U}} \end{tikzcd}}
i.e. $\Delta\vee=\deleted{\unsure{\Delta}}(\vee\otimes\vee)\Delta$, i.e $\vee$ est un
isomorphisme de co-algèbre.

% archive sheet number: 165
\scanpage{167}
\uline{Compatibilité combinée de $\varepsilon,\vee,\Delta$}, On veut que
pour $V\in K^{\underline U}$, l'hom naturel $V'\otimes V\longrightarrow k$
soit un $\underline U$-hom, quand $\underline k$ est considéré comme
$U$-module trivial (si on n'a pas supposé $\Delta$
symétrique, il faut supposer la $=$ chose pour
$V\otimes V'\longrightarrow k$). Ceci s'exprime par
\[ \underline{P(\vee\otimes1)\Delta=1}\ \text{et}\ \underline{P(\deleted{\unsure{(}}1\otimes\vee)\Delta=1} \]
(ces deux sont équivalentes si $\Delta$ est symétrique)

Si \uline{tous} ces axiomes sont vérifiés, alors
tout hom. fonctoriel classique construit sur
des $\underline U$-modules permis est permis.

\deleted{On peut \illegible{3} \unsure{qui soit} \unsure{tout}}

\noindent\rule{2.5cm}{0.4pt}

Soient $\underline U,\underline V$ deux ``$K$-algèbres'' avec
hom diagonal $\Delta$, $\varphi$ un hom \inserted{\uline{continu}} $\underline V\longrightarrow\underline U$
i.e. $\overline\varphi$ un foncteur permis $K^{\underline V}\longrightarrow K^{\underline U}$.
\uline{Pour que $\overline\varphi$ soit compatible avec les}
\uline{produits tensoriels} (i.e. si $V,W\in K^{\underline U}$,
alors sur $V\otimes W$ la structure de $\underline V$-module
associée est le produit \unsure{des st|deux st} à $\underline\varphi(\unsure{W|V})$ et $\underline\varphi(\unsure{w|W})$) il faut et il suffit
que $\varphi$ soit compatible avec $\Delta$.

% archive sheet number: 166
\scanpage{168}
Si \uline{toutes} les données et tous les axiomes sont
réunies sur $\underline U$, on a ce qui suit ($k$ un \unsure{corps comm.}).
\[
\left\{\vcenter{\hbox{\begin{minipage}{0.88\linewidth}
$\underline U$ algèbre associative topologique complète, avec suffisamment de représentations continues de dim. finie

$\underline U$ muni d'une \deleted{\unsure{homom.|structure} d'algèbre \unsure{cont} $\underline U\xrightarrow{\Delta}\underline U\otimes\underline U$}\inserted{structure de co-algèbre associative} \deleted{\unsure{iden}}\inserted{et} commutative

Les deux structures compatibles\,: \unsure{c'est|$\Delta$ est} un hom. continu d'algèbre associative

$\underline U$ muni d'une augmentation $\varepsilon$ qui est un hom continu d'algèbre et un hom. de Co-Algèbre

$\underline U$ muni d'une involution $\vee$, qui est un antiautomorphisme \inserted{continu} d'algèbre associative, et un automorphisme de coalgèbre;

$\varepsilon$ \unsure{invariante|invariant} sous $\vee$\,: $\varepsilon(\breve x)=\varepsilon(x)$

On a \quad $P(\vee\otimes1)\Delta=P(1\otimes\vee)\Delta=1$
\end{minipage}}}\right.
\]
($P$ désigne en principe l'hom $\underline U\mathbin{\widehat\otimes}\underline U\longrightarrow\underline U$
qui correspond au produit, \deleted{pour que ce soit}\inserted{il est bien}
\deleted{défini, il faudrait que pour tout $V\in K^{\underline U}$,}
\deleted{existe un $W\in K^{\underline U}$ tel que $\mathcal J_W U\subset$}
défini par passage au quotient).

\uline{Par dualité}, ceci correspond aux
données suivantes sur un ensemble $A$
\[
\left\{\vcenter{\hbox{\begin{minipage}{0.85\linewidth}
\deleted{\unsure{$A$}}: exactement les $=$s, sauf qu'il n'y a plus de topologie, et que c'est pour la structure d'algèbre ass. et non de co-algèbre que $A$ est \unsure{commutative|cocommutative}
\end{minipage}}}\right.
\]
% archive sheet number: 167
\scanpage{169}
3) Le groupe $G$ des éléments inversibles de $U$
tels que \uuline{$\Delta s=s\otimes s$,} \inserted{$\varepsilon(s)=1$. Ces $s$} \deleted{\unsure{Ils}} correspondent aux
\deleted{\unsure{V}}automorphismes du foncteur $I^{\underline U}$ tels que,
si $A$ et $B$ \deleted{sont} correspondent à $V,W\in K^{\underline U}$,
alors $A\otimes B$ correspond à $V\otimes W$. [Ils
\margin{La représentation de $G$ dans $V\otimes W$ est le produit tensoriel de ses représentations dans $V$ et $W$}
% slanted
correspondent aussi aux homomorphismes
\deleted{dans $\underline U$} continus, compatibles avec $\Delta$,
dans $\underline U$, de l'algèbre $\underline U(\mathrm Z)$ correspondant
au groupe additif des entiers et les
représentations \deleted{de \unsure{groupes}}\inserted{linéaires} de dim finie
de ce groupe]. On conclut de
\margin{$G$ opère \unsure{trivialemt} dans le module trivial. La représentation de \deleted{$G$ opérant sur $V$} $G$ dans $V'$ est contragréd. de la représentation de $G$ dans $V$}
% slanted
$\Delta s=s\otimes s$ que $\varepsilon(s)=\varepsilon(s)^2$, donc $\varepsilon(s)$ est égal
à $0$ ou $1$, mais pas $0$, car $s$ est supposé
inversible, d'où \uline{$s\in G\Longrightarrow\varepsilon(s)=1$}.

De plus, la \deleted{\unsure{de}} condition $P(\unsure{w|\vee|v}\otimes1)\Delta=1$
prouve $\breve s s=1$, d'où \uline{$\breve s=s^{-1}$}.

Conclusion\,: \uline{Si $G$ est total dans $U$}, il
ne peut exister qu'une seule application
\deleted{\unsure{continue} $\unsure{\Rightarrow}$} linéaire $\vee:U\longrightarrow U$, continue,
telle que $P(\vee\otimes1)\Delta=1$. Cette $\vee$ est
automatiquement un anti-automorphisme
d'algèbre associative, satisfait à $\varepsilon(\breve x)=\varepsilon(x)$
et est \deleted{compatible avec la structure}\inserted{un automorphisme de Co-algèbre}.
Elle est aussi caractérisée par $s^\vee=s^{-1}$
si $s\in G$. \deleted{\unsure{D'où}}

$\Delta s=s\otimes s$ signifie que $f\longrightarrow\langle s,f\rangle$
\deleted{\unsure{continue}} est un hom. de l'algèbre
commutative $A=\underline U'$ (dual

% archive sheet number: 168
\scanpage{170}
N.B. Les \unsure{élémts|éléments} de $A$ (formes linéaires continues
sur $U$) sont exactement les formes du type
$u\longrightarrow\langle u.a,b'\rangle$, où $a,b'\in V\in K^{\underline U}$, i.e. des
\unsure{coéfficients|coëfficients|coefficients} de représentation linéaires continues
de dim finie de $\underline U$. \deleted{\unsure{Ces représent.}}

\begin{center}\shortstack{\rule{1.4cm}{0.4pt}\\[-0.9ex]\rule{1.2cm}{0.4pt}}\end{center}

\deleted{Éléments}
\uline{Classes remarquables d'éléments dans $\underline U$}

1) Le groupe $U^*$\deleted{\unsure{$G$}} des éléments inversibles de $U$.
Correspond aux \uline{automorphismes} du foncteur
$I^{\underline U}$, \deleted{qui} i.e. la \deleted{\unsure{donnée pour tout}}\inserted{donnée pour} tout $\unsure{V|W}\in K^{\underline U}$
d'un \deleted{\unsure{automor|automorphisme}} $k$-automorphisme de $V$,
avec condition de commutation pour les
$\underline U$-homomorphismes. \uline{Il y en a beaucoup}
\uline{trop}.

\unsure{3|2}) Le \deleted{\unsure{groupe des élém|groupe des éléments}}\inserted{sous-espace $\unsure{\mathfrak g|g}$ des} éléments primitifs
Il s'agit des endomorphismes du foncteur
additif $I^{\underline U}$, tel que si $A$ et $B$ sont
les opérations dans $V$ et $W$, celle dans
$V\otimes W$ soit $A\otimes1+1\otimes B$. $\unsure{\mathfrak g|g}$ est une
algèbre de Lie \inserted{topologique complète}, (et $=$ une $p$-algèbre de
Lie si on est en car. $p\ne0$)
\margin{voir plus bas une étude plus \unsure{complète}}
% slanted

% archive sheet number: 169
